%!TEX root = ../report.tex
\documentclass[report.tex]{subfiles}
\begin{document}

\chapter{Evaluation and Results}
\label{ch:results}

The four sections follow the four experiments of Section~\ref{sec:timing}:
choose the backbone, choose the flow backend, characterise the coverage
threshold, then run the full validation split.

\section{Segmentation Backbone}
\label{sec:backbones}

The encoders were trained behind the same UPerNet decoder under the recipe of
Table~\ref{tab:trainhp}; Table~\ref{tab:backbones} compares them.

\begin{table}[b]
    \centering
    \caption[Trained segmentation models]{Trained segmentation models.
    Validation mIoU is at each model's best epoch on the $1024^2$ centre-crop
    validation transform (Section~\ref{sec:metric}).}
    \label{tab:backbones}
    \footnotesize
    \begin{tabular}{@{}lrrr@{}}
        \toprule
        \textbf{Model} & \textbf{Params} & \textbf{Inference} & \textbf{Val mIoU} \\
        \midrule
        ConvNeXt-Large & \SI{234.9}{M} & \LatConvnext~ms       & \textbf{\MIoUConvnextLarge} \\
        Swin-Large     & \SI{233.9}{M} & \LatSwin~ms           & \MIoUSwinLarge \\
        Hiera-Base+    & \SI{101.0}{M} & \LatHieraBasePlus~ms  & \MIoUHieraBasePlus \\
        \bottomrule
    \end{tabular}

    \vspace{1mm}
    \footnotesize\raggedright
    Inference is a single-frame \SI{4}{K} measurement on the A100. That run wrote no
    persistent log, so unlike every other number in this chapter these three
    latencies come from project notes.
\end{table}

ConvNeXt-Large has both the highest mIoU and the lowest latency, so the
pipeline uses it in every subsequent experiment. Swin-Large is within
\num{0.009} mIoU but \SI{68}{ms} slower per frame, and Hiera-Base+ trades
\num{0.036} mIoU for a \SI{57}{\percent} smaller parameter count
(Figure~\ref{fig:segqual}).

\begin{figure}[t]
    \centering
    \includegraphics[width=\linewidth]{fig_deck_segmentation.jpg}
    \caption[ConvNeXt-Large segmentation on a validation frame]{ConvNeXt-Large
    on seq36 frame 100: buildings in red, road in purple, trees and vegetation
    in green. The frame scores \num{0.917} mIoU against its dense ground
    truth.}
    \label{fig:segqual}
\end{figure}

\section{Optical Flow Backend}
\label{sec:backends}

SEA-RAFT and FlowNet2 were compared on \SeaRaftPairs{} matched
keyframe-to-target propagation pairs scored against dense ground truth
(Table~\ref{tab:backends}).

\begin{table}[b]
    \centering
    \caption[Flow backend comparison]{SEA-RAFT against FlowNet2 on
    \SeaRaftPairs{} matched Ruralscapes propagation pairs. Propagation mIoU is
    the all-pixel figure, the valid-pixel row excludes pixels the FB check
    rejects, and coverage is the mean fraction of pixels it accepts.}
    \label{tab:backends}
    \footnotesize
    \begin{tabular}{@{}lrr@{}}
        \toprule
         & \textbf{SEA-RAFT} (default) & \textbf{FlowNet2} \\
        \midrule
        Inference per pair (\SI{4}{K}) & \SI{150}{ms}            & \SI{532}{ms} \\
        Throughput                     & 6.7 pairs/s             & 1.9 pairs/s \\
        Propagation mIoU               & \SeaRaftMIoUAll         & \FlowNetTwoMIoUAll \\
        mIoU on valid pixels           & \SeaRaftMIoUValid       & \FlowNetTwoMIoUValid \\
        Valid coverage                 & \SeaRaftCoverage\%      & \FlowNetTwoCoverage\% \\
        \bottomrule
    \end{tabular}

    \vspace{1mm}
    \footnotesize\raggedright
    The per-pair timings are development measurements whose log was not
    retained; the quality and coverage figures are re-derived from the
    committed per-pair metrics files.
\end{table}

The backends are indistinguishable on quality: propagation mIoU differs by
\num{0.0006}, valid-pixel mIoU by \num{0.0001} and coverage by
\SI{0.4}{\percent}. SEA-RAFT is $3.6\times$ faster per pair and simpler to
deploy, so the choice rests on speed and maintainability rather than accuracy.
The result also confirms that \texttt{scale}${=}{-}2$
(Section~\ref{sec:flowbackends}) costs under \num{0.01} mIoU while removing
more than two thirds of the flow time.

\paragraph{Why this experiment uses Ruralscapes.} It inherited a single-jump
propagation harness built for that dataset and was not re-run on UAVid. The
regime therefore differs from the pipeline's, one flow estimate spanning 50
frames rather than 50 chained consecutive estimates, and Ruralscapes has
twelve classes against UAVid's six, so these absolute mIoU values are not
comparable with Section~\ref{sec:valresults}. The experiment is used only for
the \emph{relative} ranking of two backends on identical pairs.

\section{What the Coverage Threshold Controls}
\label{sec:sweep}

The threshold $\tau$ was swept over $\{0.75, 0.90, 0.95, 1.00\}$ on seq16 with
ConvNeXt-Large and SEA-RAFT held fixed (Table~\ref{tab:sweep}).

\begin{table}[t]
    \centering
    \caption[Coverage threshold sweep]{Threshold sweep on seq16, 901 frames,
    ConvNeXt-Large with SEA-RAFT. $\tau = 1.00$ degenerates to the
    model-every-frame baseline, since coverage can never exceed 1.}
    \label{tab:sweep}
    \footnotesize
    \begin{tabular}{@{}lrrr@{}}
        \toprule
        \textbf{Threshold $\tau$} & \textbf{Model calls}
        & \textbf{Speed vs.\ every frame} & \textbf{mIoU} \\
        \midrule
        0.75            &   5 & $2.2\times$ & 0.51 \\
        0.90            &  23 & $2.1\times$ & 0.62 \\
        0.95 (default)  &  71 & $1.9\times$ & 0.71 \\
        1.00 (baseline) & 901 & $1.0\times$ & 0.76 \\
        \bottomrule
    \end{tabular}

    \vspace{1mm}
    \footnotesize\raggedright
    The $\tau = 0.75$ and $\tau = 0.90$ rows are development runs whose
    per-frame logs were not retained.
\end{table}

\paragraph{The threshold controls quality, not speed.} This follows from the
cost structure of Figure~\ref{fig:system}. Dropping from 901 model calls to 71
removes \SI{92}{\percent} of the inference work but yields only $1.9\times$,
because flow runs on every frame regardless; going further, to 5 calls, moves
only to $2.2\times$ while costing \num{0.20} mIoU. Equivalently, raising
$\tau$ from 0.75 to 1.00 lifts mIoU from \num{0.51} to \num{0.76} for under
\SI{15}{\percent} more runtime. There is therefore no reason to operate at a
low threshold, and $\tau = 0.95$ is used throughout.

\section{Validation Split Results}
\label{sec:valresults}

Baseline and hybrid were run over all \NumSeq{} UAVid validation sequences,
\NumFrames{} frames, at $\tau = 0.95$ (Table~\ref{tab:valresults}).

\begin{table}[!htb]
    \centering
    \caption[Baseline against adaptive hybrid on the validation
    split]{Baseline against adaptive hybrid on all \NumSeq{} UAVid validation
    sequences, $\tau = 0.95$, ConvNeXt-Large with SEA-RAFT. Times are
    wall-clock over all 901 frames of each sequence; mIoU is the mean over that
    sequence's ten dense ground-truth frames, scored on identical pixel sets
    for both methods.}
    \label{tab:valresults}
    \footnotesize
    \begin{tabular}{@{}lrrrrrr@{}}
        \toprule
        \textbf{Sequence} & \textbf{Model calls} & \textbf{Baseline}
        & \textbf{Hybrid} & \textbf{Speedup}
        & \textbf{Base mIoU} & \textbf{Hyb.\ mIoU} \\
        \midrule
        seq16 &  71 & \SI{337}{s} & \SI{181}{s} & $1.86\times$ & 0.763 & 0.711 \\
        seq17 & 208 & \SI{334}{s} & \SI{210}{s} & $1.59\times$ & 0.780 & 0.758 \\
        seq18 &  76 & \SI{334}{s} & \SI{184}{s} & $1.82\times$ & 0.835 & 0.813 \\
        seq19 & 194 & \SI{364}{s} & \SI{222}{s} & $1.64\times$ & 0.744 & 0.728 \\
        seq20 & 474 & \SI{363}{s} & \SI{279}{s} & $1.30\times$ & 0.773 & 0.739 \\
        seq36 & 131 & \SI{365}{s} & \SI{211}{s} & $1.73\times$ & 0.842 & 0.831 \\
        seq37 &  87 & \SI{335}{s} & \SI{189}{s} & $1.77\times$ & 0.759 & 0.708 \\
        \midrule
        \textbf{All} & \textbf{\MeanCalls} & \textbf{\SI{\BaseTime}{s}}
        & \textbf{\SI{\HybTime}{s}} & $\mathbf{\Speedup\times}$
        & \textbf{0.785} & \textbf{0.755} \\
        \bottomrule
    \end{tabular}
\end{table}

The hybrid is \Speedup$\times$ faster than model-every-frame inference,
retains \Retention\% of the baseline mIoU (\HybMIoU{} against \BaseMIoU{}),
and calls the model on \PctModel\% of frames, \TotalCalls{} calls across
\NumFrames{} frames.

\paragraph{Model calls follow motion, without per-sequence tuning.} Under one
fixed threshold the high-motion seq20 uses 474 model calls and the calm seq37
uses 87. At that extreme the speedup falls to $1.30\times$, the lowest in the
table, but quality holds at \num{0.739} against a \num{0.773} baseline: the
cost of motion is paid in speed rather than accuracy. On the propagated frame
of Figure~\ref{fig:hybqual} the gap to the baseline is \num{0.005} mIoU.

\begin{figure}[!htb]
    \centering
    \includegraphics[width=\linewidth]{fig_deck_hybrid.jpg}
    \caption[Baseline against adaptive hybrid on a propagated frame]{Left:
    ConvNeXt-Large on every frame. Right: the adaptive hybrid at
    $\tau = 0.95$ on the same frame, labelled by propagation rather than
    inference (\texttt{src=flow}) at \SI{97}{\percent} coverage. The mIoU
    difference is \num{0.005}.}
    \label{fig:hybqual}
\end{figure}

\section{Where Quality Is Lost}
\label{sec:whereloss}

Propagation costs \num{0.030} mIoU in aggregate, from \BaseMIoU{} to
\HybMIoU, concentrated on small moving objects, principally \texttt{person}
and \texttt{vehicle}. Three properties combine to produce this. First, mIoU
averages over classes without weighting by pixel count
(Section~\ref{sec:metric}), so a \texttt{person} region of a few hundred
pixels loses a large \emph{fraction} of its IoU to a boundary ring a few
pixels wide. Second, every step resamples with nearest-neighbour
interpolation, adding a rounding error at every boundary that accumulates
along a run, invisible to the consistency check because the flow fields agree
while the boundary drifts. Third, independently moving objects cause the
occlusion that breaks correspondence. A rejected pixel keeps its propagated
label regardless: the check decides \emph{when to re-infer}, not what to do
with one bad pixel.

\section{Threats to Validity}
\label{sec:threats}

\begin{itemize}
    \item \textbf{Quality rests on \NumGT{} evaluation points.} Quality is
          measured at ten frames per sequence while timing covers all 901. The
          comparison is paired, so the \emph{difference} is well founded, but
          the sample is small. Early stopping also selected each checkpoint on
          these frames, so absolute mIoU is optimistic.
    \item \textbf{Single dataset, one altitude regime, one seed.} All pipeline
          results are UAVid; the flow-backend comparison is Ruralscapes, in a
          different propagation regime. Each backbone was trained once, so
          Table~\ref{tab:backbones} supports a ranking but not a significance
          claim.
    \item \textbf{One hardware configuration.} Every timing comes from the
          same A100, and the speedup depends on the ratio of model cost to
          flow cost, so it would differ on an embedded accelerator.
\end{itemize}

\end{document}
